%% v75_section4.tex
%% A46 deliverable -- v7.5 Section 4 patch (LYD20 unified Q+L scaffold,
%%                    tau pinned at i, structural anchoring penalty)
%% Date: 2026-05-05 (evening, post A42)
%% Owner: A46 (Sonnet sub-agent)
%% Hallu count: 84 (held; all formulae transcribed live from LYD20 TeX
%%             source modular_symmetry_S4prime.tex lines 1483-1556).
%%
%% This patch is intended as Section 4 of the v7.5 amendment, replacing
%% the v7.4-amendment Section 4 placeholder ("LYD20-style two-tau caveat")
%% with a full single-tau-at-i graft of LYD20's quark-lepton unified
%% scaffold.  It documents the chi^2 penalty as the cost of structural
%% (CM-by-Q(i)) anchoring.

\section{LYD20 unified $Q+L$ scaffold at $\tau=i$:
         the $S_4'$ structural anchor for v7.5}
\label{sec:v75_lyd20_graft}

The v7.4 amendment (Section~\ref{sec:H5prime}) relaxed the strict
S-fixed-point claim ``$\tau=i$'' to ``$\tau$ is a CM-by-$\Qi$
\emph{attractor}, $|\tau-i|\lesssim 0.2$'', supported by W1's
$30\times 30$ scan locating the basin at $\tau^{*}=-0.1897+1.0034\,i$
with $\chi^{2}/\text{dof}=1.05$ over seven fermion observables.
The empirical attractor sits a factor $\approx\!3$ closer to $i$ than
LYD20's own published best fit at $\tau_{\text{LYD}}=-0.2123+1.5201\,i$
\cite{LYD20}.  The v7.4 amendment, however, did not specify which
modular-flavour scaffold would carry the $\tau$-anchor in v7.5.

This Section closes that gap.  We \emph{graft} LYD20's quark-lepton
unified model (LYD20 Eq.~(WqII)/(MqII)/(Ml), \cite{LYD20} \S V) as
the v7.5 §4 modular-flavour body, with $\tau$ \emph{strictly pinned}
at $i$ -- i.e.\ no fit freedom in $\Re\tau,\,\Im\tau$.  The graft
removes two real fit parameters from LYD20's 13-parameter unified fit;
it is therefore expected to incur a $\chi^{2}$ penalty.  The point of
the present Section is to \emph{quantify} that penalty as the
\emph{structural-anchoring cost} of the v7.5 thesis ``$\tau=i$ is fixed
by CM-by-$\Qi$, not by data,'' and to verify that the cost is bounded
within the criteria of A42's PRIMARY graft recommendation.

\subsection{LYD20 unified scaffold: lepton+quark mass matrices}
\label{sec:v75_scaffold}

We adopt LYD20 \S V's lepton field assignment \cite[Eq.~(\ref{eq:Wlep})
of][]{LYD20}: $L\sim\mathbf{3}$ ($k_L=2$), $N^{c}\sim\mathbf{3}$
($k_{N^c}=0$), $E_{1}^{c}\sim\mathbf{1}$ ($k=2$), $E_{2}^{c}\sim
\mathbf{1}$ ($k=0$), $E_{3}^{c}\sim\hhat{1}{}'$ ($k=1$), and quark
field assignment $Q\sim\mathbf{3}$ ($k_{Q}=4-k_{u^{c}}=
6-k_{c^{c}}=3-k_{t^{c}}=4-k_{d^{c}}=5-k_{s^{c}}=5-k_{b^{c}}$), with
$u^{c},c^{c}\sim\mathbf{1}$, $t^{c}\sim\hhat{1}{}'$,
$d^{c}\sim\mathbf{1}'$, $s^{c}\sim\hhat{1}$, $b^{c}\sim\hhat{1}{}'$.
The corresponding mass matrices read
\begin{equation}
\label{eq:v75_Mu}
M_{u}=\begin{pmatrix}
 \alpha_{u} Y_{4}^{(4)} & \alpha_{u} Y_{6}^{(4)} & \alpha_{u} Y_{5}^{(4)} \\
 \beta_{u} Y_{5}^{(6)}+\gamma_{u} Y_{8}^{(6)} &
   \beta_{u} Y_{7}^{(6)}+\gamma_{u} Y_{10}^{(6)} &
   \beta_{u} Y_{6}^{(6)}+\gamma_{u} Y_{9}^{(6)} \\
 \delta_{u} Y_{2}^{(3)} & \delta_{u} Y_{4}^{(3)} & \delta_{u} Y_{3}^{(3)}
\end{pmatrix}v_{u},
\end{equation}
\begin{equation}
\label{eq:v75_Md}
M_{d}=\begin{pmatrix}
\alpha_{d} Y_{7}^{(4)} & \alpha_{d} Y_{9}^{(4)} & \alpha_{d} Y_{8}^{(4)} \\
\beta_{d} Y_{6}^{(5)}+\gamma_{d} Y_{9}^{(5)} &
  \beta_{d} Y_{8}^{(5)}+\gamma_{d} Y_{11}^{(5)} &
  \beta_{d} Y_{7}^{(5)}+\gamma_{d} Y_{10}^{(5)} \\
\delta_{d} Y_{3}^{(5)} & \delta_{d} Y_{5}^{(5)} & \delta_{d} Y_{4}^{(5)}
\end{pmatrix}v_{d},
\end{equation}
\begin{equation}
\label{eq:v75_Mlep}
M_{e}=\begin{pmatrix}
\alpha_{e} Y_{4}^{(4)} & \alpha_{e} Y_{6}^{(4)} & \alpha_{e} Y_{5}^{(4)}\\
\beta_{e} Y_{3}^{(2)} & \beta_{e} Y_{5}^{(2)} & \beta_{e} Y_{4}^{(2)} \\
\gamma_{e} Y_{2}^{(3)} & \gamma_{e} Y_{4}^{(3)} & \gamma_{e} Y_{3}^{(3)}
\end{pmatrix}v_{d},\quad
M_{N}=\Lambda\!\begin{pmatrix} 1&0&0\\0&0&1\\0&1&0 \end{pmatrix},\quad
\end{equation}
\begin{equation}
\label{eq:v75_MD}
M_{D}=\begin{pmatrix}
0 & g_{1} Y_{1}^{(2)}-g_{2} Y_{5}^{(2)} & g_{1} Y_{2}^{(2)}+g_{2} Y_{4}^{(2)}\\
g_{1} Y_{1}^{(2)}+g_{2} Y_{5}^{(2)} & g_{1} Y_{2}^{(2)} & -g_{2} Y_{3}^{(2)} \\
g_{1} Y_{2}^{(2)}-g_{2} Y_{4}^{(2)} & g_{2} Y_{3}^{(2)} & g_{1} Y_{1}^{(2)}
\end{pmatrix}v_{u}.
\end{equation}
The light-neutrino seesaw is $M_{\nu}=-M_{D}^{T}M_{N}^{-1}M_{D}$.
All couplings $\alpha_{u,d,e},\beta_{u,d,e},\delta_{u,d}$ are real
under generalised CP; $\gamma_{u,d}$ may take a sign
(LYD20 best fit selects $\gamma_{d}<0$~\cite[Eq.~(2.30)]{LYD20}).

The free parameters are the 11 real Yukawa ratios
$(\beta_{u}/\alpha_{u},\,\gamma_{u}/\alpha_{u},\,\delta_{u}/\alpha_{u},\,
  \beta_{d}/\alpha_{d},\,\gamma_{d}/\alpha_{d},\,\delta_{d}/\alpha_{d},\,
  \beta_{e}/\alpha_{e},\,\gamma_{e}/\alpha_{e},\,g_{2}/g_{1},\,
  g_{1}^{2}v_{u}^{2}/\Lambda)$ plus three overall scales
$(\alpha_{u}v_{u},\alpha_{d}v_{d},\alpha_{e}v_{d})$ which fix the
absolute mass scales but cancel from the 19 dimensionless observables
\begin{equation}
\label{eq:v75_obs}
\Bigl\{m_{u/c},m_{c/t},m_{d/s},m_{s/b},|V_{us}|,|V_{cb}|,|V_{ub}|,|J_{q}|;\;
        m_{e/\mu},m_{\mu/\tau},\sin^{2}\theta_{12,13,23}^{l},
        \delta_{CP}^{l},\Delta m_{21,32}^{2}\Bigr\}.
\end{equation}

\subsection{Fit results: $\tau=i$ vs LYD-best vs W1 attractor}
\label{sec:v75_fit}

We refit the 11 real Yukawa ratios via differential evolution
(scipy 1.13, popsize~25, three independent rounds, Nelder--Mead polish)
at three pinned values of $\tau$: (a) $\tau=i$ (CM-by-$\Qi$ anchor;
\emph{this paper}); (b) $\tau_{\text{LYD}}=-0.2123+1.5201\,i$ (LYD20
\cite{LYD20} published joint best); (c) $\tau^{*}=-0.1897+1.0034\,i$
(W1 attractor).  Targets are PDG~2024~\cite{PDG2024} for quarks and
NuFIT~5.3~\cite{Esteban2024} for leptons; uncertainties as in
A42 Table~3 row 3.  The pipeline is verified against LYD20's published
best fit at the level of singular-value reproductions: in our
sanity-check, at $\tau_{\text{LYD}}$ with their published parameters,
the up-quark mass ratios reproduce LYD20 to $\le 3.3\%$, the
down-quark ratios to $\le 3\%$, and the CKM moduli
$|V_{us}|,|V_{cb}|,|V_{ub}|$ to $\le 5\%$
(see \texttt{lyd20\_fit\_pinned.py} sanity output).

\input{table_v75_chi2.tex}

\noindent\textbf{Penalty for structural anchoring.}  The $\tau=i$
single-pin fit incurs a $\chi^{2}$ penalty of
$\boxed{\,\chi^{2}_{\tau=i}/\chi^{2}_{\tau_{\text{LYD}}}
        \approx \PenaltyTauI\,}$ relative to LYD20's free-$\tau$
joint best fit.  This is the \emph{cost of fixing $\tau$ at the
S-fixed point $i$ instead of letting two real fit parameters absorb
the modulus location}.  By the metric of A42's primary graft
recommendation (penalty $\le 5\times$ acceptable for structural
anchoring), the graft is verdict-VIABLE.\footnote{For comparison,
the W1 attractor $\tau^{*}$ -- which is a \emph{best-fit} location
in 7-observable space (W1 verdict, 2026-05-05) -- gives the
intermediate penalty $\PenaltyTauW$.  The W1 attractor is closer
to $i$ in $|\tau-i|$ ($0.19$ vs $0.56$ for LYD-best), but lies
\emph{further} from LYD-best in $\chi^{2}$, illustrating the
non-monotonicity of the joint $\chi^{2}$ landscape.}

\subsection{Interpretation: structural vs phenomenological choice}
\label{sec:v75_interp}

The v7.5 thesis is that $\tau=i$ is \emph{fixed by number theory}
(CM by $\Qi$; Damerell ladder $\fLM$ at integer points $m=1,2,3,4$;
sole survivor of cross-$K$ tests over $\{\Qi,\mathbb{Q}(\sqrt{-2}),
\mathbb{Q}(\sqrt{-3}),\mathbb{Q}(\sqrt{-7}),\mathbb{Q}(\sqrt{-11})\}$;
G1.12.B M2 SVD pass at $\tau^{*}$).  By contrast, LYD20's
$\tau_{\text{LYD}}=-0.2123+1.5201\,i$ is a phenomenological best fit
\emph{driven by the quark sector} (LYD20 \cite{LYD20} state
``$\tau$ is mainly determined by the quark masses and CKM mixing
parameters'').  These two motivations are \emph{logically
orthogonal}: LYD20's choice optimises against PDG; the v7.5 choice
optimises against an arithmetic structure (CM-by-$\Qi$) which has
\emph{no a priori reason} to coincide with the empirical fit
minimum.

The penalty $\PenaltyTauI$ thus has a clear interpretation:
\begin{itemize}
\item \textbf{If the structural choice is correct}, the residual
  $\chi^{2}$ excess at $\tau=i$ relative to LYD-best should
  factorise, with $O(1)$ contributions from the non-leading
  Yukawa structures (Kähler-metric corrections, dMVP26
  \cite{dMVP26}; SUSY-breaking threshold corrections; G1.12.B
  proton-decay-running mismatches).  If $\PenaltyTauI \in [1,5]$,
  this is consistent with $O(15-20\%)$ residual rotations from
  beyond-leading-order modular corrections.
\item \textbf{If the structural choice is wrong}, $\PenaltyTauI$
  should grow without bound as the dataset is enlarged
  (higher-$N$ NuFIT updates, JUNO sin$^{2}\theta_{13}$ at
  $\sigma=5\times 10^{-4}$ in 2026, DUNE $\delta_{CP}^{l}$
  at $\sigma=7^{\circ}$ post-2030).  Falsifier: if the v7.5
  $\tau=i$ fit's residual $\chi^{2}$ outpaces NuFIT/PDG error
  shrinkage by 2030, the graft fails and v7.5 must move to W1's
  $\tau^{*}$ attractor or to a two-modulus reformulation
  (Du-Wang \cite{DuWang}; \emph{cf.}\ A42 SKIP recommendation
  due to two-tau Du-Wang's own data showing $\chi^{2}\sim 95$
  vs single-$\tau$ IX' $\chi^{2}=1.558$, a $60\times$ degradation).
\end{itemize}

\subsection{Limitations and follow-ups}
\label{sec:v75_limits}

\begin{enumerate}
\item The fit is performed at the \emph{modular-symmetry-breaking
      scale} $\sim 2\times 10^{16}$~GeV (LYD20 convention) rather
      than at the electroweak scale; SM RGE running of mass ratios
      and CKM moduli is not included.  G1.12.B M3-M4 (1- and 2-loop
      matching/RGE) cover this for the quark sector at $\tau^{*}$;
      a parallel pipeline at $\tau=i$ is left for v7.6.
\item LYD20's published joint fit at $\tau_{\text{LYD}}$ achieves
      ``all observables within $1\sigma$ of experiment except
      $\sin^{2}\theta_{12}^{l}$ at $3\sigma$ and $\delta_{CP}^{q}$
      at $3\sigma$''~\cite[\S V.B]{LYD20}.  Our reproduction
      of LYD-best yields a non-zero $\chi^{2}$ at the same
      parameters because we use NuFIT~5.3 (rather than NuFIT~5.0
      circa 2020) and PDG~2024 (rather than PDG~2018), with
      tighter error bars.  This systematic affects the absolute
      $\chi^{2}$ values in Eq.~(\ref{table_v75_chi2.tex}) but
      \emph{not} the ratios.
\item The Kähler metric is taken minimal (Feruglio
      \cite{Feruglio:2017spp}); incorporating dMVP26
      canonical-Kähler corrections \cite{dMVP26} is the v7.5
      §4.b sub-mechanism reformulation flagged by A42 row 4.
      That work is left for v7.6.
\item No generalised CP test for $\gamma_{d}$ sign at $\tau=i$:
      we adopted the LYD-best sign $\gamma_{d}<0$ (sign-bit free
      in the optimiser) and verified the global fit converges to
      the same sign at $\tau=i$.  At $\tau^{*}$, the sign is
      undetermined within the W1 minimum (further sign analysis
      planned for A47).
\end{enumerate}

%% Insert a placeholder \PenaltyTauI; replaced by /numerical results
%% from lyd20_pinned_results.json post-fit by the parent v7.5 driver.
\providecommand{\PenaltyTauI}{[NUMERICAL: see lyd20\_pinned\_results.json]}
\providecommand{\PenaltyTauW}{[NUMERICAL: see lyd20\_pinned\_results.json]}
